\section{Goldblatt-Thomason Theorem for Polyhedra}

Recall 
\begin{theorem}[Goldblatt-Thomason]
    A class of frames $C$ is definable by a set of modal formulas if and only if
    \begin{enumerate}
        \item It's closed under generated subframes
        \item It's closed under bounded morphic images
        \item It's closed under disjoint unions 
        \item It's closed under and reflects ultrafilter extensions.
    \end{enumerate}    
\end{theorem}

We want to define what it means for a class of Polyhedra to be modally definable:

\begin{definition}\hfill

    Let $C$ be a class of simplicial complexes. $C$ is modally definable if and only if the class
    \[\F(C)=\set{\Face(\Sigma)\in\Pos: \Sigma\in C}\] is modally definable
\end{definition}

And of course we need a further geometric description:
\begin{definition}\hfill 

    Let $K$ be a class of Polyhedra (seen as combinatorial manifolds), we call the \textbf{simplicial companion} of $K$ the class 
    \[\Simp(K)=\set{\Sigma\in\Scplx: |\Sigma|\in K}\]
\end{definition}
Thus we will say 
\begin{lemma}\hfill 

    A class of Polyhedra is modally definable if and ony if its simplicial companion is modally definable 
\end{lemma}
\begin{proof}
    It's shown in %\cite{} maybe logics of polyhedral reachability?
    that $|\Sigma|,x\models \phi\iff \Sigma,\sigma_x\models\phi$
    Therefore \[|-|:\set{\Sigma:\Sigma\models\set{\phi_i}_{i\in I}}\twoheadrightarrow\set{|\Sigma|:|\Sigma|\models\set{\phi_i}_{i\in I}}\]
    is an epimorphism

    We showed that $\ker(|-|)$ is the 'common triangulation' equivalence, therefore 
    \[\set{\Sigma:\Sigma\models\set{\phi_i}_{i\in I}}/\sim_{\star}\simeq \set{|\Sigma|:|\Sigma|\models\set{\phi_i}_{i\in I}}\]

    Note that $\sim_{\star}$ is always a bisimilarity, thus the quotient is modally definable if and only if the class is 'modally definable up to bisimulation'
    ,if and only if it's modally definable.
\end{proof}

Now let's see how the frame constructions are interpreted in simplicial complexes. 
\subsection{Generated subframes}

Recall:
\begin{definition}[Generated subframe]\hfill 

    Let $(K,R)$ be a Kripke frame. a \textbf{generated subframe} is a subset $K'\subseteq K$ with a relation $R'$ such that 
    \begin{itemize}
        \item if $w\in K'$ and $wRv$ then $v\in K'$.
        \item $R'=R\cap (W'\times W')$
    \end{itemize}
\end{definition}

In the face poset $R=\prec$ is the face relation. In other words we have that
\begin{proposition}\hfill 

Let $\Sigma$ be the face poset of a simplicial complex 
$\Sigma'\subseteq \Sigma$ is a generated submodel if and only if it's upwards-closed, if and only if it's bisimilar to 
\end{proposition}
\begin{proof}
For this, we regard every face poset as $\wedge-$complete, i.e., containing $\emptyset$.

Let $\Sigma$ be a simplicial complex and $\Sigma'$ a generated submodel of $\Sigma$ regarded as its face poset. 

Let $\mu$ be a minimal face in $\Sigma'$; $\mu\in \Sigma'$ implies that all faces containing $\mu$ are in $\Sigma'$ as it's upwards-closed.

Now, the case $\mu=\emptyset$ means that $\Sigma=\Sigma'$ (it's the union of all links)

The case where $\mu$ is a maximal face is also trivial: the only singleton-poset is the poset only containing $\emptyset$, which is the empty union.

If $\mu$ is neither minimal nor maximal, then consider two cases: Either $\mu$ is the only minimal element or it isn't.

If $\mu$ is the only minimal element, thus $\Sigma'=\uparrow\mu$: Every face of $\uparrow\mu$ contains $\mu$; Remove $\mu$ from every element of $\uparrow\mu$
and call $\ell$ the resulting poset. 
\[\ell=\set{\sigma\setminus\mu:\sigma\in\uparrow\mu}\]
Note that now every element of $\ell$ is an element of $\Sigma$ as $\Sigma$ is downward-closed, and $\sigma\cap \mu=\emptyset$ by definition. 
This means that $\ell=\ell_{\Sigma}(\mu)$.

Now suppose there exist $\mu,\mu'$ distinct minimal elements of $\Sigma'$.
Consider the poset \[\ell=\dfrac{\uparrow\mu\cup\uparrow\mu'}{(\mu\sim\mu'\sim \emptyset)}\]

We can check that $\uparrow\mu\cup\uparrow\mu'\to \ell$ is a p-morphism:

Let $\sigma\preceq\sigma'$, then $[\sigma]\preceq[\sigma']$: The edge cases are those where $\sigma\in \uparrow\mu$ and $\sigma'=\sigma\cup\mu'$, in which cases the
quotient sends both to the same face.

Now let $[\sigma]\preceq \tau$. This means that at least one between $\tau\cup\mu$, $\tau\cup\mu'$ is in $\uparrow\mu\cup\uparrow\mu'$.

    WLOG suppose it's the first case: This means that if $[\sigma]\preceq\tau$ then in particular $\sigma\cup\mu\in \Sigma'$, thus $\tau\cup\mu$ realizes the 'zag' condition.

Now, $\ell$ is isomorphic -as a poset- to $\ell_\Sigma(\mu)\cup\ell_\Sigma(\mu')$ 
\end{proof}